% =====================================================================
%  DISSERTAÇÃO DE MESTRADO — arquivo de partida
%  Modelo ABNT do ICT/Unifei — segue a NBR 14724:2024.
%
%  Copie este arquivo, renomeie e escreva seu trabalho aqui mesmo: todas
%  as divisões do texto ficam neste único arquivo. Se preferir dividir o
%  texto em vários arquivos, veja o manual (manual.pdf) — ele próprio é
%  organizado assim e mostra como fazer.
%
%  Compile com:  latexmk modelo-dissertacao.tex
% =====================================================================
\documentclass[12pt,a4paper,oneside]{book}

% O tipo governa o modelo inteiro: natureza, folha de aprovação, resumos e
% metadados exigidos. Dissertação e tese produzem o mesmo layout, mas
% diferem no grau pretendido — para doutorado, use modelo-tese.tex.
\usepackage[tipo=dissertacao]{UnifeiICTReport}

\usepackage{csquotes}
\usepackage{fontspec}

% --- Referências bibliográficas --------------------------------------
% Para citar trabalhos e produzir a lista de referências, crie o SEU
% arquivo .bib (por exemplo, referencias.bib, nesta mesma pasta) e
% descomente as duas linhas abaixo e o \printbibliography lá embaixo.
%
% ATENÇÃO: não aponte para o referencias.bib deste repositório. Ele é
% fictício — traz referências inventadas que servem apenas aos exemplos
% do manual. Um trabalho entregue com ele tem bibliografia falsa, e isso
% não salta aos olhos ao folhear o PDF.
%
%\usepackage[style=abnt]{biblatex}
%\bibliography{referencias.bib}

% --- Siglas e símbolos ------------------------------------------------
% Declare aqui as siglas do seu trabalho. Depois de declarar ao menos
% uma, descomente também o \printglossary correspondente lá embaixo.
% Sem nenhuma declaração, a lista sairia como uma página de título
% seguida de nada.
%
%\makeglossaries
%\newacronym{ict}{ICT}{Instituto de Ciências e Tecnologia}

%----------------------------------------------------------------------
% Dados da capa e da folha de rosto
%----------------------------------------------------------------------

\title{Título da sua dissertação}

% Escreva só o texto do subtítulo, sem os dois-pontos: a NBR 14724:2024
% §4.1.1(d) exige que ele venha "precedido de dois-pontos", e o modelo
% insere essa pontuação sozinho. Comente a linha se não houver subtítulo.
%\subtitle{Subtítulo, se houver}

\author{Seu Nome Completo}

\supervisor{Prof. Dr. Nome do Orientador}

%\cosupervisor{Prof. Dra. Nome da Coorientadora}

% Programa de pós-graduação. Entra na natureza do trabalho.
\subject{Nome do Programa de Pós-Graduação}

% OBRIGATÓRIA em dissertação e tese: a NBR 14724:2024 §4.2.1.1.1(e)
% inclui a área de concentração na natureza do trabalho, e o modelo
% interrompe a compilação se ela faltar.
\areaconcentracao{Área de Concentração}

%\linhapesquisa{Linha de Pesquisa}

% --- Banca examinadora ------------------------------------------------
% Declare apenas os membros EXTERNOS: o orientador já entra pela linha
% \supervisor acima e não deve ser repetido aqui. Campos em branco ainda
% produzem um bloco corretamente formatado, para a folha preparada antes
% de os nomes serem conhecidos.
\bancamembro{Prof. Dr. Nome do Membro}{Doutor em Área}{Unifei}
\bancamembro{Prof. Dra. Nome da Membra}{Doutora em Área}{Instituição}

%\aprovacaodata{15 de dezembro de 2026}

%----------------------------------------------------------------------
% Resumos
%----------------------------------------------------------------------

\abstract{%
    Escreva aqui o resumo da sua dissertação: objetivos, metodologia,
    resultados e conclusões, em um único parágrafo.
}
\keywords{primeira palavra-chave; segunda; terceira.}

% OBRIGATÓRIO neste tipo de trabalho: a NBR 14724:2024 §4.2.1.8 torna o
% resumo em língua estrangeira elemento obrigatório, e o modelo
% interrompe a compilação se ele faltar.
\abstractseclang{%
    Write the abstract in the foreign language here.
}
\keywordsseclang{first keyword; second; third.}

%======================================================================
\begin{document}

\frontmatter

\maketitle

\folhaaprovacao

\makeabstracts

% --- Listas -----------------------------------------------------------
% Descomente cada lista DEPOIS de o trabalho ter o que listar. Chamada
% sem conteúdo produz uma página com o título e nada abaixo.
%\listoffigures
%\listoftables
%\listofquadros
%\listofgraficos
%\printglossary[type=\acronymtype,title={\acronymlistname}]
%\printglossary[type=symbols,title={\symbolslistname}]

\tableofcontents

\mainmatter

% Atenção ao vocabulário: você digita \chapter, e o documento imprime
% "Seção". A NBR 14724:2024 §4.2.2 não admite dividir o trabalho em
% capítulos, então o modelo rebaixa cada nível um degrau:
%   \chapter -> Seção          \subsection    -> Subsubseção
%   \section -> Subseção       \subsubsection -> Parágrafo
% O manual explica isso em detalhe.

\chapter{Introdução}\label{cap:introducao}

Contextualize o tema, delimite o problema, apresente os objetivos e
descreva como o trabalho está organizado.

\chapter{Referencial teórico}\label{cap:referencial}

Apresente os conceitos e trabalhos que fundamentam a pesquisa.

\chapter{Metodologia}\label{cap:metodologia}

Descreva como a pesquisa foi conduzida.

\chapter{Resultados e discussão}\label{cap:resultados}

Apresente os resultados obtidos e discuta-os à luz do referencial.

\chapter{Conclusão}\label{cap:conclusao}

Retome os objetivos, destaque as contribuições e aponte trabalhos
futuros.

% Descomente junto com o bloco de bibliografia do preâmbulo.
%\printbibliography

\backmatter

% Apêndices são de sua autoria; anexos são documentos de terceiros.
% Descomente o grupo que for usar (veja o manual para os detalhes).
%\apendices
%\apendice{ape:rotulo}{Título do apêndice}
%
%\anexos
%\anexo{ane:rotulo}{Título do anexo}

\end{document}
